\chapter{Altered Mental Status}

\section*{Definition}
Altered mental status is any change in attention, awareness, thinking, behavior, or level of consciousness that comes on over hours to days. It is a symptom -- not a disease -- and it signals that something is disturbing your brain. The most common forms are delirium (a fluctuating confusion, most often in older adults) and a reduced level of arousal from drowsiness to deep unresponsiveness. Because the cause can be a stroke, low blood sugar, an infection, or a poisoning, a sudden change in mental status is always an emergency.

\section*{When to See a Doctor}
Call for emergency help if:
\begin{itemize}
  \item A person suddenly cannot be awakened, is hard to rouse, or stops answering sensibly
  \item The change is sudden or the person has a seizure, drooping face, weakness of one side, or slurred speech
  \item Confusion is accompanied by a fever and a stiff neck (possible meningitis)
  \item The person has diabetes and low blood sugar, or is on blood-thinning medicines and has fallen
  \item There is vomiting, head injury, poisoning, or any rapid deterioration over hours
\end{itemize}
Treat this like a medical emergency -- do not wait to see if it passes.

\section*{How It Develops}
Your brain depends on a constant supply of oxygen and glucose and on a stable chemical environment. In delirium, many disorders temporarily disrupt the brain's balance of neurotransmitters such as acetylcholine and dopamine, so nerve cells lose their normal coordination and the brain struggles to focus attention. In a depressed level of consciousness, the problem is often widespread -- both sides of the brain or the structures that keep you awake (the reticular activating system) are affected by low oxygen, low blood sugar, drugs, toxins, liver or kidney failure, or a very high or low body temperature. A stroke or head injury can depress consciousness by directly damaging the arousal centers. What looks like confusion can also be a seizure, a severe mood disorder, or a metabolic derangement such as a very high calcium level.

\section*{Causes}
\begin{itemize}
  \item Infections: urinary tract, chest, or blood infections; meningitis and encephalitis; and sepsis
  \item Metabolic derangements: low blood sugar, low oxygen, dehydration, very low or high sodium, calcium, or ammonia, and liver or kidney failure
  \item Stroke, bleeding in the brain, or transient ischemic attack
  \item Head injury, including concussion
  \item Medications and substances: sedatives, opioids, anticholinergics, and steroids; alcohol intoxication or withdrawal
  \item Seizures, including subtle or non-convulsive seizures
  \item Organ failure and cardiorespiratory problems: heart attack, heart failure, and chronic lung disease
  \item In older people, even a mild infection or a medication change can trigger delirium
\end{itemize}

\section*{Related Symptoms}
\begin{itemize}
  \item Difficulty paying attention or following a conversation
  \item Disorientation to time, place, or people
  \item Wandering, restlessness, or agitation; conversely, sleepiness and withdrawal
  \item Hallucinations, usually visual, and paranoia
  \item Slurred speech, tremors, or jerking movements
  \item Fever, stiff neck, headache, or light sensitivity
  \item Cold, clammy skin or sweating with confusion (low blood sugar)
\end{itemize}

\section*{Medical Evaluation}
\begin{itemize}
  \item A bedside check of alertness, orientation, and attention using a tool such as the Glasgow Coma Scale or the Confusion Assessment Method.
  \item Immediate blood tests for glucose, oxygen, electrolytes, and kidney and liver function.
  \item A careful medication and substance history, since drug effects and withdrawal are leading causes.
  \item Blood tests and cultures looking for infection, plus a head CT scan when stroke, bleeding, or injury is possible.
  \item An ECG, and, if a seizure is suspected, an EEG; a lumbar puncture may be needed if meningitis is a concern.
\end{itemize}

\section*{Treatment Options}
\begin{itemize}
  \item Correct the underlying trigger: glucose for hypoglycemia, oxygen for hypoxia, fluids and antibiotics for infection, and antidotes or supportive care for poisons.
  \item Stop or reduce offending medicines, especially sedatives and anticholinergics.
  \item Treat seizures with anti-seizure medicines when a seizure is the cause.
  \item For agitated delirium, brief use of low-dose antipsychotics such as haloperidol may be needed while the cause is treated.
  \item Non-drug care matters most: keeping lights and routines normal, involving family, reorienting the person, and preventing falls and pressure sores.
\end{itemize}


\section*{Self-Care and Home Management}
\begin{itemize}
  \item For a person who is confused but otherwise stable, keep them calm, oriented (remind them of the date and place), and safe from wandering.
  \item Do not give food, drink, or oral medicine to anyone who is drowsy or unable to swallow -- they could choke.
  \item If diabetes is present and the person can swallow safely, give sugar only if you can confirm or strongly suspect low blood sugar; otherwise, seek help.
  \item Gather the medicine list, any substances taken, and the time the change began for the emergency team.
  \item Never try to "sleep off" a sudden change in mental status -- get medical attention.
\end{itemize}

\section*{References}
\begin{thebibliography}{99}
\bibitem{altered_mental_status:delirium} Inouye SK. Delirium in older persons. \textit{N Engl J Med}. 2006;354:1157--1165.
\bibitem{altered_mental_status:dsm} American Psychiatric Association. \textit{Diagnostic and Statistical Manual of Mental Disorders}. 5th ed., text revision. American Psychiatric Association; 2022.
\bibitem{altered_mental_status:ed-eval} Kanich W, Brady WJ, Huff JS, et al. Altered mental status: evaluation and etiology in the emergency department. \textit{Am J Emerg Med}. 2002;20:613--617.
\bibitem{altered_mental_status:nomenclature} Slooter AJC, Otte WM, Devlin JW, et al. Updated nomenclature of delirium and acute encephalopathy: joint task force. \textit{Intensive Care Med}. 2020;46:1020--1022.
\end{thebibliography}
